\ifdefined\XeTeXversion
  % IEEEtran selects ptm before loading the user preamble. Map its Unicode
  % NFSS families explicitly so XeTeX retains serifs, italics and bold.
  \newcommand{\UXfont}[5]{\DeclareFontShape{TU}{#1}{#2}{#3}{<->"[#4]:mapping=tex-text"}{}}
  \newcommand{\UXfamily}[5]{%
    \DeclareFontFamily{TU}{#1}{}%
    \UXfont{#1}{m}{n}{#2}{}\UXfont{#1}{m}{it}{#3}{}%
    \UXfont{#1}{b}{n}{#4}{}\UXfont{#1}{bx}{n}{#4}{}%
    \UXfont{#1}{b}{it}{#5}{}\UXfont{#1}{bx}{it}{#5}{}%
    \DeclareFontShape{TU}{#1}{m}{sl}{<->ssub*#1/m/it}{}%
    \DeclareFontShape{TU}{#1}{m}{sc}{<->ssub*#1/m/n}{}}
  \IfFileExists{times.ttf}{%
    \UXfamily{ptm}{times.ttf}{timesi.ttf}{timesbd.ttf}{timesbi.ttf}
    \UXfamily{phv}{arial.ttf}{ariali.ttf}{arialbd.ttf}{arialbi.ttf}
    \UXfamily{pcr}{cour.ttf}{couri.ttf}{courbd.ttf}{courbi.ttf}
  }{%
    \UXfamily{ptm}{texgyretermes-regular.otf}{texgyretermes-italic.otf}{texgyretermes-bold.otf}{texgyretermes-bolditalic.otf}
    \UXfamily{phv}{texgyreheros-regular.otf}{texgyreheros-italic.otf}{texgyreheros-bold.otf}{texgyreheros-bolditalic.otf}
    \UXfamily{pcr}{texgyrecursor-regular.otf}{texgyrecursor-italic.otf}{texgyrecursor-bold.otf}{texgyrecursor-bolditalic.otf}
  }
\fi
\documentclass[conference]{IEEEtran}
\usepackage[brazilian]{babel}
\usepackage{array,tabularx,amsmath,graphicx}
\newcommand{\institution}{CEFET-MG, Campus Divinópolis, Engenharia de Computação}
\newcommand{\groupauthors}{\IEEEauthorblockN{Élcio Costa Amorim Neto, Guilherme Alvarenga de Azevedo e Joaquim Cézar Santana da Cruz}\IEEEauthorblockA{\institution}}
\addto\captionsbrazilian{\renewcommand{\abstractname}{Resumo}\renewcommand{\refname}{Referências}}
\renewcommand{\IEEEkeywordsname}{Palavras-chave}
\newcommand{\UXev}[1]{\textsf{\detokenize{#1}}}
\newcommand{\UXrow}[2]{\noindent\begin{tabularx}{\columnwidth}{@{}>{\raggedright\arraybackslash}p{.30\columnwidth}>{\raggedright\arraybackslash}X@{}}\textbf{#1} & #2\end{tabularx}\par}
\newenvironment{uxtable}[2]{\par\smallskip\begingroup\footnotesize\setlength{\parskip}{1pt}\setlength{\tabcolsep}{3pt}\hrule\nobreak\UXrow{#1}{\textbf{#2}}\nobreak\hrule\nobreak}{\nobreak\hrule\endgroup\par\smallskip}

\AtBeginDocument{\raggedbottom}


\title{Instrumentação de um cardápio móvel para auditar escolhas e conferir pedidos sob ampliação}
\author{\groupauthors}
\begin{document}
\maketitle
\begin{abstract}
Uma pessoa que usa a ampliação de tela pode ver o preço de uma massa sem ver o nome do prato ao qual ele se refere. Uma observação textual que reúna essas duas informações anteciparia uma associação que a pessoa precisa estabelecer. Este artigo descreve a implementação de um cardápio instrumentado para investigar essa diferença de acesso à informação. As duas versões preservam trinta alimentos, seus preços e as operações essenciais da comanda. A primeira organiza os itens em lista contínua. A segunda oferece categorias e mantém um resumo do pedido no rodapé. O controlador recorta os signos pela geometria da interface, limita a memória de trechos vistos e determina os gestos disponíveis antes de obter a escolha do agente. A trilha de eventos distingue exposição, proposta, ação executada e estado do pedido. Um verificador externo compara a comanda confirmada com o objetivo da tarefa. Quatro execuções determinísticas verificaram a instrumentação em duas tarefas de calibração, nas duas versões. Para pedir a lasanha individual de menor preço, o método especialista consumiu oito ações na lista e sete na versão por categorias. Para corrigir a quantidade de uma sobremesa no pedido iniciado pela irmã, consumiu três ações em cada versão. As quatro comandas atenderam aos critérios de sucesso e a contagem de ações coincidiu com a referência independente. O pacote reúne ainda nove sessões Gemini do comparador C: uma teve sucesso, duas atingiram o limite de ações e seis terminaram por falha técnica. As tarefas de validação e a comparação com participantes permanecem reservadas para investigação posterior.
\end{abstract}
\begin{IEEEkeywords}acessibilidade, baixa visão, engenharia semiótica, modelos de interação, avaliação de usabilidade\end{IEEEkeywords}

\section{Introdução e pergunta de projeto}
Se uma comanda confirmada contém uma sobremesa a menos, o pesquisador precisa distinguir três fatos: quais linhas entraram na vista, quais operações alteraram a sobremesa e o que foi registrado ao confirmar. A mensagem de confirmação só estabelece o último fato. Observar a interface pelo texto integral da página encobriria a dificuldade que motivou o estudo de cardápios sob ampliação \cite{fase1,pesquisa}.

A pergunta desta fase é quais decisões sobre a interface, a ajuda e a conexão com um modelo de linguagem permitem auditar uma tarefa com a informação que uma persona teria. Tratamos a interface como a mensagem do projetista, recebida tanto pela pessoa quanto por um agente \cite[p.~84]{desouza}. O menu fornece os mesmos signos de domínio aos dois receptores. A representação textual e os mecanismos externos especificam a aproximação que se pode testar. O caderno E1 especifica signos e falas.

\section{Modelagem herdada}
Lívia é assistente administrativa de 34 anos, personagem fictícia com baixa visão congênita, campo visual preservado e leitura fluente quando a palavra cabe na vista. Usa ampliação de tela cheia e fonte grande no Android. O estudo delimita adultos entre 18 e 59 anos que recorrem à ampliação e dispensam leitor de tela. A hierarquia de objetivos, operações e planos segue a análise hierárquica de tarefas \cite{annett}. Os modelos anteriores do grupo definem o cardápio, os estados iniciais, o comparador e a separação entre calibração e validação \cite[apêndices B--D]{modelagem}\cite[pp.~4--9]{hts}.

C1 requer comparar três lasanhas individuais e pedir uma unidade da de menor preço. C2 requer conferir quatro linhas e adequar a comanda ao combinado. O programa inicia C2 no cardápio com a comanda carregada e registra o acesso à revisão, conciliando a narrativa do pedido iniciado com o plano de três toques da HTA. V1 reserva o jantar vegetariano para dois sob limite de R\$\,110,00. V2 reserva a repetição de um jantar com alteração de tamanho e retirada de um item residual. O piloto só aceita C1 e C2.

O perfil P sorteia ampliação entre 3 e 5, adota limiar de legibilidade 3 e varredura linear. P2 e P6 variam os intervalos para 1,5--2,5 e 5,5--7 e os limiares para 1,5 e 3. Nessas variantes, a probabilidade de estratégia linear é 9/11. P, P2, P6 e C compartilham capacidade de três a cinco elementos lembrados e o mesmo acionador de ajuda. C inicia em ampliação unitária com varredura livre e probabilidade de deslize zero. G recebe a cena corrente inteira como controle genérico. R executa regras determinísticas.

\section{Decisões e estatuto epistemológico}
A alternância entre contexto e alvo durante uma busca ampliada foi observada em computador \cite[p.~9]{billah}. Em smartphone, a literatura descreve diferenças de estratégia e incidentes gestuais \cite[pp.~12, 20, 25]{islam}. A transferência para preços, opções e linhas de um restaurante é uma hipótese operacional. A Tabela~\ref{tab:dec} distingue as decisões e suas bases.
\begin{table}[ht]
\caption{Decisões que afetam a interpretação dos registros}
\label{tab:dec}
\footnotesize\centering
\begin{tabularx}{\columnwidth}{@{}>{\raggedright\arraybackslash}p{.22\columnwidth}>{\raggedright\arraybackslash}X>{\raggedright\arraybackslash}p{.22\columnwidth}@{}}
\hline Decisão & Justificativa & Estatuto\\ \hline
Variar ampliação e varredura & Os estudos de origem registram diferenças no acesso a alvos e contexto. & Evidência indireta\\
Manter nome e atributo como signos separados & Permite registrar qual deles efetivamente entrou na vista. & Princípio de Design\\
Usar o mesmo $k,n,q$ em P e C & A curadoria não fundamenta atribuir menor memória nem outra disposição para pedir ajuda ao perfil. & Hipótese Operacional\\
Permitir rever e editar & A prevenção por agrupamento mantém uma via para recuperar seleções equivocadas. & Princípio de Design\\
Fixar resumo em B & A composição à vista pode reduzir retornos, mas também cobrir alimentos e sair da ampliação. & Hipótese Operacional\\
Compartilhar ajuda e guardas & Os receptores precisam ter acesso à mesma explicação e às mesmas operações de domínio. & Princípio de Design\\ \hline
\end{tabularx}
\end{table}

\section{Arquitetura e acoplamento}
A aplicação serve as opções admissíveis de acordo com a cena e a composição. Em S2, uma massa exige a escolha de tamanho antes da inclusão. A resposta a uma alteração é a mesma no navegador e no executor. A transição entre as cenas segue o diálogo da MoLIC \cite{molic}. As categorias de signos distinguem a expressão física de sua interpretação \cite[p.~230]{barbosa}. O modelo físico deriva do conceitual conforme a distinção de Benyon \cite[p.~199]{benyon}.

\begin{figure}[ht]
\centering
\setlength{\unitlength}{.99pt}
\begin{picture}(238,136)
\put(3,106){\framebox(102,26){\shortstack{Aplicação\\estado e geometria}}}
\put(131,106){\framebox(103,26){\shortstack{M1--M3, M5 e M6\\recortar e restringir}}}
\put(105,119){\vector(1,0){26}}
\put(183,106){\vector(0,-1){20}}
\put(131,60){\framebox(103,26){\shortstack{Política\\escolher ação}}}
\put(3,60){\framebox(102,26){\shortstack{M4, M7 e guardas\\aplicar ação válida}}}
\put(131,73){\vector(-1,0){26}}
\put(54,86){\vector(0,1){20}}
\put(54,60){\vector(0,-1){27}}
\put(3,7){\framebox(102,26){\shortstack{Registro do analista\\eventos e estado final}}}
\put(105,20){\vector(1,0){26}}
\put(131,7){\framebox(103,26){\shortstack{Verificador externo\\resultado e métricas}}}
\end{picture}
\caption{Ciclo de execução e leitura independente dos resultados. O acionador de ajuda M7 também pode solicitar uma ação por impasse. O estado integral alimenta a auditoria.}
\label{fig:arch}
\end{figure}

A Figura~\ref{fig:arch} mostra os dois usos do estado. Para preparar a próxima decisão, M1 calcula a interseção de cada signo com uma janela de $360/z$ por $800/z$ pixels CSS. Textos parcialmente expostos são truncados. Abaixo do limiar, restam pistas geométricas e o tipo de elemento. M2 conserva fragmentos que saíram da vista, em ordem de chegada e com capacidade $k$. M3 omite a associação ao nome quando o cartão aparece em partes. M5 aplica o recorte também na revisão. M6 restringe a navegação conforme a estratégia. A pessoa pode solicitar ajuda a qualquer instante pelo acesso ubíquo. As solicitações explícitas e as ações de protocolo permanecem distinguíveis de toques em signos expostos.

O adaptador serializa a orientação funcional, um exemplo do treino com pão de queijo, a instrução literal e a projeção atual. A resposta deve conter uma única escolha válida. O guardião compara a projeção serializada com a produzida pelo canal e rejeita chaves analíticas ou estruturas proibidas. Após uma resposta malformada, o sistema permite duas novas tentativas. M4 pode transformar um deslocamento em um toque no ponto inicial, com probabilidade 0,05 nos perfis P. M7 observa cinco ações consecutivas sem progresso e pode solicitar ajuda com probabilidade 0,5. A interface registra essa abertura antes de expor o conteúdo de S5. A projeção, as decisões e os dados integrais do pesquisador seguem caminhos distintos.

\section{Instrumentação e rupturas observáveis}
Cada evento associa sessão, configuração, tarefa, versão, semente, passo, instante, cena, região visível e origem sintética. O resumo SHA-256 identifica a configuração amostrada. Os adaptadores arquivam o corpo de cada tentativa, o modelo solicitado, a resposta e, quando disponível, sua identificação efetiva e latência. O vocabulário fechado de 23 eventos discrimina exposição de signo, proposta, rejeição, execução, mutação e encerramento. Uma consulta à ajuda identifica a camada e a origem. As seis condições de suporte incluem disponibilidade e discrição \cite[pp.~397--399]{dix}. Um mesmo texto em três camadas é usado pelos dois receptores.

Operacionalizamos seis etiquetas do Método de Avaliação de Comunicabilidade \cite{mac}. Abrir ajuda gera \emph{Socorro!}; desistir explicitamente gera \emph{Desisto}. Confirmar um pedido divergente produz \emph{Pra mim está bom}. Um deslize seguido em até duas ações por retorno ou desfazer produz \emph{Epa!}. Uma sequência de cinco ações sem novo signo, mudança de cena ou mutação do pedido produz \emph{Cadê?}, uma vez por sequência. O acesso à dica local sobre um tópico produz \emph{O que é isso?}. Essas correspondências são hipóteses operacionais. O mesmo evento pode satisfazer duas regras.

Para a omissão em C1, a exposição requer o valor individual completo. Em C2, o verificador reúne por linha um prefixo legível do nome (até quatro caracteres) e os dígitos da quantidade. Essa aproximação da identificação é uma hipótese operacional. O verificador recebe o estado final e os eventos. Em C2, examina todas as linhas da refeição e a confirmação, inclusive o acompanhamento do frango e as duas unidades de pudim. $S$ conta a execução iniciada e bem-sucedida entre as sessões elegíveis. Falha técnica fica fora do denominador. $Re_{E_k}$ divide erros pelas oportunidades da categoria. E1 corresponde à seleção, E2 à associação, E3 à quantidade ou opção, E4 ao deslize, E5 à divergência inserida que persistiu e E6 às restrições de V1. Uma categoria sem oportunidade recebe valor indefinido. $A$ conta ações executadas e $R_A=A/N_{\rm ref}(z)$. $D$ mede deslocamentos e mudanças de ampliação por alvo localizado, $V$ mede revisitas, $O$ mede omissões antes da decisão e $U$ conta acessos por camada de ajuda.

\section{Piloto e referência especialista}
O pacote E3 contém uma sessão de R para cada combinação de C1/C2 e A/B, com semente 42 e $z=1$. Para cada sessão, um segundo percurso calcula $N_{\rm ref}$ na mesma geometria. A Tabela~\ref{tab:pilot} apresenta o resultado do recálculo do verificador sobre os registros arquivados. O custo hipotético por escolha independente segue a Tabela I do KLM \cite[p.~399, Tabela~I]{card80}: preparação $M=1{,}35$~s, apontamento $P=1{,}10$~s e acionamento $K=0{,}20$~s. A aproximação de um toque por $PK$ pressupõe alvo previamente localizado. O KLM se aplica ao especialista sem erros \cite[pp.~448--456]{preece}.

\begin{table}[ht]
\caption{Percursos de R em ampliação unitária, uma sessão por linha}
\label{tab:pilot}\centering\small
\begin{tabular}{@{}lrrrrrrr@{}}\hline
Tarefa/versão & Toques & Desl. & $A$ & $N_{\rm ref}$ & $S$ & $R_A$ & $U$\\ \hline
C1/A & 5 & 3 & 8 & 8 & 1 & 1 & 0\\
C1/B & 6 & 1 & 7 & 7 & 1 & 1 & 0\\
C2/A & 3 & 0 & 3 & 3 & 1 & 1 & 0\\
C2/B & 3 & 0 & 3 & 3 & 1 & 1 & 0\\ \hline
\end{tabular}
\end{table}
Em C1, a categoria de B introduziu um toque e o percurso exigiu dois deslocamentos a menos. Em C2, abrir a revisão, corrigir a quarta linha e confirmar produziram as três ações previstas nas duas versões. Sob a premissa dos três alvos localizados, a sequência $3(MPK)$ corresponde a 7,95~s de operadores, sem modelar leitura, espera ou deslizes. Em C1 os deslocamentos observados impedem interpretar o produto entre toques e tempo por escolha como duração total. R registrou zero ocorrências em todas as categorias de erro. As taxas definidas foram zero e as oportunidades nulas permaneceram indefinidas. As quatro execuções tiveram $O=0$. Em C1/A, $D=1$ e $V=0$; em C1/B, $D=1/3$ e $V=1/3$. Em C2, ambos os índices foram zero. A proporção de sessões com ajuda foi zero em cada camada.

O E3 consolidado contém 13 registros: quatro de R e nove do comparador C com Gemini 3.5 Flash Lite, semente 42. Em 486 tentativas foram arquivadas 478 respostas. O protocolo previa 16 condições, das quais três alcançaram desfecho comportamental. As seis falhas técnicas, cinco em C1/A e uma em C2/B, são excluídas do denominador de sucesso. C2/B executou 64 ações antes da interrupção técnica. C1/A e C1/B atingiram o limite de 150 ações, com $S=0$, $U=0$, respectivamente um erro E1 e 47 erros E3. Em C2/A, $S=1$, $A=N_{\rm ref}=3$, $R_A=1$ e $U=0$. Em C2/A, $Re_{E3}=0/1$ e $Re_{E5}=0/1$. As demais taxas ficam indefinidas por ausência de oportunidades. Os logs, denominadores e taxas por tentativa estão no resumo de E3. O visor permite acompanhar as três decisões da sessão \texttt{C2-A-C-42-840b102d725e}. O instrumento de coleta corresponde à revisão \texttt{fc0ba7a} do repositório do grupo. O E1 acompanha este artigo como apêndice separado. Os arquivos de teste com erros conhecidos verificam os cálculos do instrumento e permanecem identificados como material sintético de teste.

\section{Ameaças à validade}
O protótipo inclui uma extensão demonstrativa de navegação manual com fotografias ilustrativas em Detalhes do item (S2). As condições A e B dos resultados deste artigo têm cardápio sem fotos. As imagens da demonstração ficam fora da observação da política e da coleta. Avaliar seu efeito exigiria versionar a geometria e definir como preservar a equivalência de informação entre a pessoa e o agente.
O recorte opera sobre caixas e uma aproximação tipográfica. Coordenadas textuais, posição de alternativas e seleção por índice diferem da percepção, do contraste e do gesto em um telefone real. A exposição de parte de um termo não comprova leitura. Memória FIFO e deslize probabilístico são escolhas do modelo operacional. O controle G acessa a cena corrente completa, condição deliberadamente distinta da paridade proposta para P e C. Modelos remotos podem alterar comportamento mesmo quando o identificador solicitado e o texto permanecem iguais.

Uma sessão por condição verifica apenas a execução do método. Os estudos de origem diferem em dispositivo, tarefa ou população. Falta uma amostra humana pareada. O protocolo anterior propõe contrastar diferenças P--C em dados humanos e sintéticos, com intervalos de 90\% e tolerâncias pré-fixadas \cite[p.~7]{hts}. Antes dessa comparação serão necessários revisão ética, dados independentes e congelamento de artefatos, parâmetros, ajuda, modelos e critérios.

\section{Conclusão}
A ligação entre geometria, fala admissível, mudança da comanda e verificação externa permitiu explicar, ação por ação, a diferença técnica de C1 entre as versões e a coincidência dos percursos de C2. Os quatro registros sustentam a reconstituição desses percursos. Os registros remotos verificam o percurso do comparador C. A avaliação da plausibilidade da persona requer execuções de P e uma referência humana independente. A revisão de coleta não simulava duração de hesitação: tempos de rede e intervalos entre requisições são operacionais. A extensão posterior permite sortear uniformemente a hesitação entre zero e um teto explícito, com semente independente, e registrar o atraso simulado por proposta. Esta hipótese permanece sem calibração humana. O protocolo preserva V1 e V2 para essa etapa.

\bibliographystyle{IEEEtran}\bibliography{referencias}
\end{document}
